\documentclass[11pt]{article}

\usepackage[T1]{fontenc}
\usepackage[utf8]{inputenc}
\usepackage{lmodern}
\usepackage{microtype}
\usepackage[margin=1in]{geometry}
\usepackage{amsmath,amssymb,amsthm,mathtools}
\usepackage{booktabs}
\usepackage{array}
\usepackage{enumitem}
\usepackage{xcolor}
\usepackage{url}
\usepackage[colorlinks=true,linkcolor=blue!55!black,citecolor=green!40!black,urlcolor=blue!60!black]{hyperref}
\usepackage[nameinlink,noabbrev]{cleveref}

\newtheorem{theorem}{Theorem}[section]
\newtheorem{lemma}[theorem]{Lemma}
\newtheorem{proposition}[theorem]{Proposition}
\newtheorem{corollary}[theorem]{Corollary}
\newtheorem{definition}[theorem]{Definition}
\newtheorem{conjecture}[theorem]{Conjecture}
\newtheorem{problem}[theorem]{Open Problem}
\theoremstyle{remark}
\newtheorem{remark}[theorem]{Remark}
\newtheorem{example}[theorem]{Example}

\newcommand{\R}{\mathbb{R}}
\newcommand{\Q}{\mathbb{Q}}
\newcommand{\N}{\mathbb{N}}
\newcommand{\ExistsR}{\exists\mathbb{R}}
\newcommand{\cR}{\mathcal{R}}
\newcommand{\cC}{\mathcal{C}}
\newcommand{\cE}{\mathcal{E}}
\newcommand{\cI}{\mathcal{I}}
\newcommand{\cP}{\mathcal{P}}
\newcommand{\Om}{\Omega}
\newcommand{\CMP}{\mathrm{CMP}}
\newcommand{\ReLU}{\operatorname{ReLU}}
\newcommand{\ExactNNTP}{\textnormal{\textsc{Exact-ReLU-NNTP-Decision}}}
\newcommand{\Feas}{\operatorname{Feas}}
\newcommand{\Proj}{\operatorname{Proj}}
\newcommand{\rowspan}{\operatorname{rowspan}}
\newcommand{\im}{\operatorname{im}}
\newcommand{\kernel}{\operatorname{ker}}
\newcommand{\tw}{\operatorname{tw}}
\newcommand{\poly}{\operatorname{poly}}
\newcommand{\bits}{\operatorname{bits}}
\newcommand{\given}{\,\middle|\,}
\newcommand{\denote}[1]{\left[\!\left[#1\right]\!\right]}

\newcommand{\proved}{\textcolor{green!45!black}{\textsc{Proved}}}
\newcommand{\conditional}{\textcolor{orange!75!black}{\textsc{Conditional}}}
\newcommand{\openstatus}{\textcolor{red!65!black}{\textsc{Open}}}
\newcommand{\refuted}{\textcolor{red!75!black}{\textsc{Refuted}}}

\setlist[itemize]{leftmargin=2em}
\setlist[enumerate]{leftmargin=2em}
\setlength{\emergencystretch}{2em}
\allowdisplaybreaks


\title{Cotangent Message Passing for Exact ReLU Feasibility\\
and the Existential Theory of the Reals}
\author{Muhammad Samuel Qudus}
\date{Draft v0.1 -- September 2026}

\begin{document}
\maketitle

\begin{abstract}
We study a decision formulation of exact training for finite feed-forward
rectified-linear-unit (ReLU) networks.  The trainable parameters are shared
across all samples, whereas preactivations, activations, and slack variables
are sample-local.  This induces a block-angular polynomial system and an
iterated tensor-product description of its coordinate ring over the shared
parameter algebra.  We establish the corresponding decomposition of relative
K\"ahler differentials, formulate an intrinsic cotangent separator message as
the image of the kernel of an internal Jacobian block, and prove its exact
factorization property for the associated linearized system.  These results
hold without selecting an activation pattern; the exact ReLU graph, including
its switching boundary, is represented by polynomial equalities with
existential square slacks.

The central question is whether these cotangent messages decide the original
zeroth-order feasibility problem.  We give an in-scope one-ReLU counterexample:
for every positive target, the evaluated scalar cotangent message is the full
separator space, even though the exact surviving parameter set is the
singleton containing that target.  Hence no decoder whose only semantic input
is the first-order message subspace can decide all exact ReLU instances.  This
separates the proved linear separator theorem from the global nonlinear
decision invariant required for a polynomial-time algorithm.  We record the
precise additional obligations---a complete zeroth-order message, a
polynomial-size representation theorem, a Turing-model bit bound, and a
model-matching reduction from the existential theory of the reals.  We also
state the resulting class collapse conditionally.  In particular, this draft
does not claim an unconditional proof that
\(P=\exists\mathbb{R}\) or \(P=NP\).
\end{abstract}

\tableofcontents

\section{Introduction}
\label{sec:introduction}

Exact feasibility questions for neural networks sit at the intersection of
real algebraic geometry, computational complexity, and formal verification.
Once approximation and numerical optimization are removed, a finite
feed-forward ReLU network with unknown real weights and rational data defines a
finite system of polynomial equalities and inequalities.  The natural decision
question is whether that system has a real solution.  Closely related neural
network training problems are known to be complete for the complexity class
\(\ExistsR\)~\cite{AbrahamsenKleistMiltzow2021}; this class contains NP and is
contained in PSPACE~\cite{Canny1988,SchaeferCardinalMiltzow2024}.

This paper investigates whether the repeated-sample structure of exact neural
training admits an exact separator algorithm.  A single parameter vector
\(\theta\) is shared across the dataset, while each sample has its own
preactivations, activations, and auxiliary variables.  At the level of
polynomial equations this yields a block-angular Jacobian.  At the level of
commutative algebra it yields an iterated tensor product over the shared
parameter algebra.  These facts motivate \emph{Cotangent Message Passing}
(CMP): eliminate a local block from the row space of a linearized constraint
system and retain only the induced subspace on a separator.

The distinction between a linearized relation and an exact existential
projection is decisive.  The cotangent module records first-order behavior.
The decision problem asks whether the semialgebraic set is empty.  Neither the
universal property of K\"ahler differentials nor a pointwise Jacobian identity
identifies these two objects.  A proof of an exact CMP decider therefore needs
an additional semantic invariant:
\begin{equation}
  s\in \mathsf{Msg}_v
  \quad\Longleftrightarrow\quad
  \exists x_{\operatorname{int}(T_v)}\;
  F_{T_v}(x,s).
  \label{eq:semantic-message-intro}
\end{equation}
It must also prove that the representation of \(\mathsf{Msg}_v\), its merge
operation, and its projection operation remain polynomially bounded.

\subsection{Contributions}

The present draft makes the following claims, each with an explicit proof
status.

\begin{enumerate}
  \item We define \(\ExactNNTP\) as a Boolean feasibility language with
  rational input and unrestricted real witnesses.  ``Training'' describes the
  origin of the constraints; the output is a Boolean, not a weight vector.

  \item We prove an exact complementarity description of a ReLU gate and its
  existential square-slack polynomialization.  The description includes the
  switching point \(z=a=0\).

  \item We derive the shared-base tensor-product description of the equality
  coordinate ring and the direct-sum decomposition of relative K\"ahler
  differentials after base change.  We then display the induced block-angular
  Jacobian.

  \item We define the intrinsic first-order CMP message
  \(
    \cE(A\mathbin{\to}S)=J_{AS}(\kernel J_A)
  \)
  and prove exact separator factorization for the associated linear row-space
  problem.  This theorem is independent of a null-space basis.

  \item We prove that the first-order message is not a complete invariant of
  exact ReLU feasibility.  A one-ReLU family with positive target \(y\) has
  exact parameter projection \(\{y\}\), but its evaluated scalar CMP message
  equals \(\R\) for every \(y>0\).  Thus targets one and two yield identical
  messages and different exact feasibility at \(\theta=1\).

  \item We isolate the remaining obligations for the proposed complexity
  collapse: a polynomially represented zeroth-order message, a correct global
  decider, polynomial Turing-model bit complexity, and a reduction landing in
  the exact language defined here.  We formalize the final class-theoretic
  implication conditionally.
\end{enumerate}

\subsection{Status of the proposed collapse}

The intended implication has the form
\begin{equation}
  \ExistsR \leq_p \ExactNNTP \in P
  \quad\Longrightarrow\quad
  P=NP=\ExistsR.
  \label{eq:collapse-chain}
\end{equation}
The implication in \cref{eq:collapse-chain} is standard.  The right-hand
premise is not supplied by the first-order CMP theorem.  Indeed,
\cref{sec:global-exactness} disproves the required completeness statement for
a decoder that receives only the existing cotangent message subspace.  The
collapse is therefore recorded as a target conditional theorem, not as an
unconditional result.

This separation is more than a matter of exposition.  The counterexample is
inside the stated problem class, uses a single ReLU, and occurs at a positive
satisfying point.  It cannot be repaired by excluding ReLU switching
singularities.  Any successful exact message must preserve zeroth-order target
constants and global component information that differentiation can erase.

\subsection{Organization}

\Cref{sec:preliminaries,sec:problem,sec:relu} define the computational and
algebraic setting.  \Cref{sec:coordinate-ring,sec:kahler,sec:block-jacobian}
establish the shared-base structure.  \Cref{sec:topology,sec:cmp} develop the
separator decomposition and the proved linear CMP theorem.
\Cref{sec:global-exactness} gives the semantic projection invariant and the
counterexample to cotangent-only completeness.  The algorithmic and complexity
consequences are audited in \cref{sec:algorithm,sec:complexity}.
\Cref{sec:etr-reduction,sec:collapse} state the exact reduction and collapse
obligations.  Boundary behavior, related work, and possible repair directions
appear in the remaining sections.

\section{Background and preliminaries}
\label{sec:preliminaries}

\subsection{The existential theory of the reals}

An existential sentence over the reals has the form
\begin{equation}
  \exists x_1,\ldots,x_m\in\R:\;
  \Psi(x_1,\ldots,x_m),
\end{equation}
where \(\Psi\) is a quantifier-free Boolean combination of polynomial sign
conditions with integer coefficients.  The decision problem ETR asks whether
such a sentence is true.  The class \(\ExistsR\) is the closure of ETR under
polynomial-time many-one reductions.  Standard containments give
\begin{equation}
  P\subseteq NP\subseteq\ExistsR\subseteq PSPACE.
\end{equation}
We use the discrete Turing model: coefficients and combinatorial structure are
encoded by finite binary strings, while quantified witnesses range over the
reals.  Surveys and algorithms for real feasibility may be found in
\cite{BasuPollackRoy2006,Renegar1992I,SchaeferCardinalMiltzow2024}; Canny's
algorithm supplies the PSPACE upper bound~\cite{Canny1988}.

\begin{definition}[Semialgebraic set]
A subset \(X\subseteq\R^d\) is semialgebraic if it is defined by a finite
Boolean combination of polynomial equalities and inequalities with real
coefficients.  A \emph{basic} semialgebraic set is defined by a conjunction of
such sign conditions.
\end{definition}

Existentially quantified square variables can replace weak nonnegativity:
\begin{equation}
  t\geq 0 \quad\Longleftrightarrow\quad \exists u\in\R:\;t=u^2.
  \label{eq:square-slack}
\end{equation}
This device permits an equality-only ambient variety whose real projection is
the intended basic semialgebraic set.  The projection operation remains part of
the semantics; the equality coordinate ring alone does not remember which
variables were introduced existentially.

\subsection{Coordinate rings and differentials}

Let \(k\) be a commutative base field and \(B\) a finitely generated
\(k\)-algebra.  The module of K\"ahler differentials \(\Om_{B/k}\) is the
\(B\)-module equipped with a universal \(k\)-derivation
\(d:B\to\Om_{B/k}\).  For a tower \(k\to A\to B\), there is a right-exact
transitivity sequence
\begin{equation}
  B\otimes_A\Om_{A/k}
  \longrightarrow \Om_{B/k}
  \longrightarrow \Om_{B/A}
  \longrightarrow 0.
  \label{eq:kahler-transitivity}
\end{equation}
These standard constructions are reviewed, for example, in
\cite{Eisenbud1995,StacksDifferentials}.

If
\begin{equation}
  B=k[x_1,\ldots,x_n]/(f_1,\ldots,f_m),
\end{equation}
then \(\Om_{B/k}\) admits the presentation
\begin{equation}
  B^m \xrightarrow{J_F^{\mathsf T}} B^n
  \longrightarrow \Om_{B/k}\longrightarrow 0,
  \label{eq:jacobian-presentation}
\end{equation}
where \(J_F=(\partial f_i/\partial x_j)\).  This is a module over the entire
quotient ring.  Evaluating at a real point \(p\) tensors the presentation with
the residue field at \(p\) and produces pointwise cotangent information.  The
symbolic module and any one evaluated vector space are distinct objects.

\subsection{Linear maps and intrinsic messages}

Let \(J_A:L_A\to X_A\) and \(J_{AS}:L_A\to X_S\) be linear maps over a field.
The intrinsic separator message is
\begin{equation}
  \cE(A\to S)
  :=J_{AS}(\kernel J_A)\subseteq X_S.
  \label{eq:intrinsic-linear-message}
\end{equation}
If the maps are represented by row blocks and \(P_A^\perp\) has rows spanning
the relevant left null space, then a matrix representing the same subspace is
\begin{equation}
  M_{A\to S}=P_A^\perp J_{AS}.
\end{equation}
The subspace, rather than the chosen basis matrix, is the invariant.

\subsection{Tree decompositions}

A tree decomposition of a graph \(G=(V,E)\) is a tree whose nodes carry bags
\(B_t\subseteq V\) such that every vertex appears in a connected family of
bags and every edge is contained in some bag.  Its width is
\(\max_t|B_t|-1\); the treewidth \(\tw(G)\) is the minimum width.  Bounded
treewidth often enables exact dynamic programming when each separator message
has a controlled representation~\cite{Bodlaender1996}.  Treewidth alone does
not bound the descriptive complexity of projected semialgebraic sets.

\subsection{Proof-status vocabulary}

To prevent a conditional dependency from being silently promoted, we use four
labels in this draft:
\begin{center}
\begin{tabular}{>{\bfseries}l p{0.72\linewidth}}
\toprule
Proved & A complete mathematical proof is present or a cited standard result
applies with its hypotheses checked.\\
Conditional & The implication is proved assuming explicitly named
certificates or lemmas.\\
Open & The statement is required by the target theorem but is not presently
derived.\\
Refuted & A counterexample disproves the statement as formulated.\\
\bottomrule
\end{tabular}
\end{center}

\section{The exact ReLU NNTP decision problem}
\label{sec:problem}

We define a finite Boolean language.  The definition is deliberately separated
from any optimization method.

\begin{definition}[Network architecture]
A feed-forward ReLU architecture is a finite directed acyclic graph with input,
affine, ReLU, and output nodes.  An affine node computes a rational affine
combination of earlier activations in which a designated subset of coefficients
are coordinates of a shared parameter vector \(\theta\in\R^P\).  Products of
an unknown weight and a sample-local activation are polynomial terms.  A ReLU
node maps a preactivation \(z\) to \(a=\max\{0,z\}\).
\end{definition}

For each sample \(i\in[N]\), write \(S^{(i)}\in\R^{K_i}\) for every
sample-local preactivation, activation, output, and auxiliary slack variable.
The parameters \(\theta\) are not copied by sample.  Rational inputs and
targets are substituted as constants.

\begin{definition}[Exact-ReLU-NNTP instance]
An instance \(I\) consists of:
\begin{enumerate}
  \item a finite feed-forward architecture and a designation of trainable
  parameters;
  \item a finite rational dataset \((x_i,y_i)_{i=1}^N\);
  \item exact affine-transition constraints, exact ReLU constraints, and exact
  output-target equalities for every sample; and
  \item a binary encoding of all rational coefficients and incidence data.
\end{enumerate}
Its feasibility predicate is
\begin{equation}
  \Feas(I);:\Longleftrightarrow\;
  \exists\theta\in\R^P\;
  \exists S^{(1)}\in\R^{K_1}\cdots
  \exists S^{(N)}\in\R^{K_N}:
  \bigwedge_{i=1}^N\Phi_i(\theta,S^{(i)}).
  \label{eq:exact-feasibility}
\end{equation}
\end{definition}

\begin{definition}[\(\ExactNNTP\)]
Given the binary encoding of \(I\), return \textnormal{\textsc{yes}} if and only if
\(\Feas(I)\) holds.
\end{definition}

No witness is required as output.  The definition does not mention a loss
surface, gradient descent, a local optimum, or numerical tolerance.  Exact
target equality is used in this draft; a threshold-loss variant is a different
formal language unless a model-matching reduction is supplied.

\subsection{Input size and uniformity}

Let \(L\) be the maximum coefficient bit length, and let \(n=|\langle
I\rangle|\) be the total binary encoding length.  The encoding explicitly
lists the architecture, trainable parameter incidences, samples, and rational
constants.  Consequently
\begin{equation}
  N\le n,\qquad P\le n,\qquad K_i\le n
  \label{eq:dimension-input-bounds}
\end{equation}
for any reasonable noncompressed encoding.  These inequalities do not by
themselves prove that an algebraic algorithm runs in polynomial time: the size
of intermediate polynomials, formulas, or algebraic numbers must also be
controlled.

\subsection{Sample locality}

The constraint family has the syntactic form
\begin{equation}
  \Phi_I(\theta,S^{(1)},\ldots,S^{(N)})
  =\bigwedge_{i=1}^{N}\Phi_i(\theta,S^{(i)}).
  \label{eq:sample-local-form}
\end{equation}
Thus distinct samples share \(\theta\) and no sample-local variable.  This
property drives both the tensor decomposition and the star-shaped dependency
graph below.

\begin{proposition}[Semantic parameter projection]
\label{prop:semantic-projection}
For each sample define
\begin{equation}
  Q_i:=\{\theta\in\R^P:\exists S^{(i)}\;Phi_i(\theta,S^{(i)})\}.
  \label{eq:sample-projection}
\end{equation}
Then
\begin{equation}
  \Feas(I)\quad\Longleftrightarrow\quad
  \bigcap_{i=1}^{N}Q_i\ne\varnothing.
  \label{eq:projection-intersection}
\end{equation}
\end{proposition}

\begin{proof}
If a common assignment \(\theta\) and local witnesses satisfy
\cref{eq:exact-feasibility}, then \(\theta\in Q_i\) for every \(i\).  Conversely,
if \(\theta\in\bigcap_iQ_i\), choose a local witness for each of the finitely
many samples.  These witnesses share no local variables and therefore combine
with \(\theta\) into a global witness.
\end{proof}

\Cref{prop:semantic-projection} is an exact zeroth-order separator theorem.  It
does not give an efficient representation of the sets \(Q_i\) or an efficient
method for testing their intersection.

\section{Exact algebraic and semialgebraic ReLU encoding}
\label{sec:relu}

The graph of a ReLU is basic semialgebraic and has a particularly small
complementarity description.

\begin{lemma}[Exact ReLU complementarity]
\label{lem:relu-complementarity}
For \(z,a\in\R\),
\begin{equation}
  a=\ReLU(z)
  \quad\Longleftrightarrow\quad
  a\ge0,\quad a-z\ge0,\quad a(a-z)=0.
  \label{eq:relu-complementarity}
\end{equation}
\end{lemma}

\begin{proof}
Suppose first that \(a=\max\{0,z\}\).  If \(z\le0\), then \(a=0\), so
\(a\ge0\), \(a-z=-z\ge0\), and the product vanishes.  If \(z\ge0\), then
\(a=z\), so both weak inequalities hold and \(a-z=0\).

Conversely, the product equation implies either \(a=0\) or \(a-z=0\).  In the
first case, \(a-z\ge0\) implies \(z\le0\), whence
\(a=0=\max\{0,z\}\).  In the second case, \(a=z\) and \(a\ge0\), whence
\(a=z=\max\{0,z\}\).
\end{proof}

The three regimes are all present:
\begin{align}
  z<0 &\Longrightarrow a=0,\\
  z=0 &\Longrightarrow a=0,\\
  z>0 &\Longrightarrow a=z.
\end{align}
No disequality \(z\ne0\) is introduced.

\begin{lemma}[Equality-only polynomialization]
\label{lem:relu-polynomialization}
For \(z,a\in\R\),
\begin{equation}
 a=\ReLU(z)
 \quad\Longleftrightarrow\quad
 \exists u,v\in\R:
 a=u^2,\quad a-z=v^2,\quad a(a-z)=0.
 \label{eq:relu-polynomialization}
\end{equation}
\end{lemma}

\begin{proof}
By \cref{lem:relu-complementarity}, a ReLU assignment makes \(a\) and \(a-z\)
nonnegative.  Each is therefore a square in \(\R\).  Conversely, the two square
equalities imply the weak inequalities, and the complementarity equation then
allows an application of \cref{lem:relu-complementarity}.
\end{proof}

\begin{remark}[Projection is part of the semantics]
The equality variety in variables \((z,a,u,v)\) is not literally the graph in
the \((z,a)\)-plane.  Its projection to \((z,a)\) is the graph.  The slack
variables therefore belong to the existential local state.
\end{remark}

\begin{proposition}[Membership in \(\ExistsR\)]
\label{prop:nntp-in-existsr}
The language \(\ExactNNTP\) belongs to \(\ExistsR\).
\end{proposition}

\begin{proof}
Introduce real variables for the shared parameters, every sample-local
preactivation and activation, and the two square slacks of each ReLU.  Affine
transitions with trainable weights are polynomial equalities of degree at most
two in the unknowns.  Exact targets are polynomial equalities.  Replace every
ReLU by \cref{eq:relu-polynomialization}.  The number of variables and
constraints and the coefficient encoding length grow polynomially with the
explicit instance encoding.  The resulting existential conjunction of
polynomial equalities is true exactly when the instance is feasible.
\end{proof}

\subsection{The switching boundary and singularity}

The complementarity polynomial
\begin{equation}
  g(z,a)=a(a-z)
\end{equation}
has gradient
\begin{equation}
  \nabla g(z,a)=(-a,,2a-z).
\end{equation}
At \((z,a)=(0,0)\), this gradient vanishes.  Thus a Jacobian-based regularity
argument cannot cover the switching point merely because the global K\"ahler
module exists there.  K\"ahler differentials are defined at singular points,
but their fibers can change rank and acquire information not present on a
smooth stratum.  The exact semantics above remains valid at the boundary; only
the first-order regularity changes.

\section{Shared-parameter coordinate-ring decomposition}
\label{sec:coordinate-ring}

We now retain the polynomial equalities, including square-slack equations, and
separate the one shared parameter algebra from sample-local algebras.  Let
\(k=\R\) for the geometric statements and put
\begin{equation}
  A=k[\theta_1,\ldots,\theta_P].
\end{equation}
For sample \(i\), let \(S^{(i)}\) denote a finite tuple of local indeterminates
and define
\begin{equation}
  B_i=A[S^{(i)}],\qquad
  \cR^{(i)}=B_i/I^{(i)},
\end{equation}
where \(I^{(i)}\) is generated by the polynomial equalities belonging to that
sample.  In the common ambient polynomial algebra
\begin{equation}
  R=A[S^{(1)},\ldots,S^{(N)}],
\end{equation}
each \(I^{(i)}\) is understood through extension from \(B_i\) to \(R\), and
\begin{equation}
  I=\sum_{i=1}^{N} I^{(i)}R,
  \qquad
  \cR=R/I.
  \label{eq:global-ring}
\end{equation}

\begin{proposition}[Ambient shared-base tensor product]
\label{prop:ambient-tensor}
There is a canonical \(A\)-algebra isomorphism
\begin{equation}
  A[S^{(1)},\ldots,S^{(N)}]
  \cong
  A[S^{(1)}]\otimes_A\cdots\otimes_A A[S^{(N)}].
  \label{eq:ambient-tensor}
\end{equation}
\end{proposition}

\begin{proof}
Both sides are the coproduct in commutative \(A\)-algebras of the polynomial
\(A\)-algebras on the disjoint variable sets \(S^{(i)}\).  Concretely, the map
sends each variable to its copy in the corresponding tensor factor.  The
inverse multiplies the images of pure tensors in the common polynomial ring.
The two maps agree on \(A\) and on every variable, and hence are inverse by the
universal property of polynomial algebras.
\end{proof}

\begin{theorem}[Quotient tensor decomposition]
\label{thm:quotient-tensor}
With the notation above, there is a canonical \(A\)-algebra isomorphism
\begin{equation}
  \boxed{
  \cR\cong
  \cR^{(1)}\otimes_A\cdots\otimes_A\cR^{(N)}.}
  \label{eq:quotient-tensor}
\end{equation}
\end{theorem}

\begin{proof}
It suffices to treat two factors and iterate.  The canonical bilinear map
\(B_1\times B_2\to (B_1\otimes_A B_2)/J\), where \(J\) is generated by the
images of \(I^{(1)}\otimes B_2\) and \(B_1\otimes I^{(2)}\), vanishes whenever
either input belongs to the corresponding ideal.  It therefore descends to a
map
\begin{equation}
  (B_1/I^{(1)})\otimes_A(B_2/I^{(2)})
  \longrightarrow (B_1\otimes_A B_2)/J.
\end{equation}
The universal property gives an inverse, so the map is an isomorphism.  Under
\cref{eq:ambient-tensor}, the ideal \(J\) corresponds to the sum of the two
extended sample ideals.  Iteration gives \cref{eq:quotient-tensor}.
\end{proof}

\begin{remark}[One parameter algebra]
The tensor product is over \(A\).  It identifies the parameter coordinates in
all local factors.  A tensor product over \(k\) would create independent copies
of the parameters and would model a different feasibility problem.
\end{remark}

\subsection{Equality rings versus feasible real sets}

\Cref{thm:quotient-tensor} concerns polynomial equalities.  The exact decision
problem concerns real points together with the existential interpretation of
slack variables.  Passing to a quotient ring does not perform existential
projection, choose a real connected component, or decide whether real points
exist.  These distinctions matter because two systems can have closely related
differential modules while having different real feasibility.

The global real feasible set may be written
\begin{equation}
  X_I(\R)=\left\{(\theta,S^{(1)},\ldots,S^{(N)}):
  F^{(i)}(\theta,S^{(i)})=0\text{ for all }i\right\},
\end{equation}
after square-slack polynomialization.  The accepted parameter set is its
projection to \(\theta\), equivalently the intersection in
\cref{eq:projection-intersection}.  The coordinate-ring decomposition explains
the sample incidence pattern, but it does not compute this projection.

\section{K\"ahler differential structure}
\label{sec:kahler}

The shared-base tensor product has a corresponding decomposition of relative
differentials.

\begin{theorem}[Relative differential decomposition]
\label{thm:relative-differentials}
Let \(\cR\) be the iterated tensor product in
\cref{eq:quotient-tensor}.  There is a canonical \(\cR\)-module isomorphism
\begin{equation}
  \boxed{
  \Om_{\cR/A}\cong
  \bigoplus_{i=1}^{N}
  \left(
    \Om_{\cR^{(i)}/A}
    \otimes_{\cR^{(i)}}\cR
  \right).}
  \label{eq:relative-differential-decomposition}
\end{equation}
\end{theorem}

\begin{proof}
For two \(A\)-algebras \(B,C\), an \(A\)-derivation from \(B\otimes_A C\) to a
\((B\otimes_A C)\)-module \(M\) is uniquely determined by its restrictions to
\(B\) and \(C\).  Conversely, two such derivations define
\(d(b\otimes c)=c\,d_B(b)+b\,d_C(c)\).  The universal property of K\"ahler
differentials therefore gives
\begin{equation}
 \Om_{(B\otimes_A C)/A}
 \cong
 (\Om_{B/A}\otimes_B(B\otimes_A C))
 \oplus
 (\Om_{C/A}\otimes_C(B\otimes_A C)).
\end{equation}
Iterating this isomorphism and applying \cref{thm:quotient-tensor} proves the
claim.
\end{proof}

This theorem is global and requires no smoothness hypothesis.  Its meaning is
relative: after the shared parameter algebra is held fixed, different samples
have independent differential generators and relations.

\begin{proposition}[Absolute-to-relative transitivity]
\label{prop:transitivity}
For \(k\to A\to\cR\), there is a right-exact sequence
\begin{equation}
  \cR\otimes_A\Om_{A/k}
  \longrightarrow\Om_{\cR/k}
  \longrightarrow\Om_{\cR/A}
  \longrightarrow0.
  \label{eq:global-transitivity}
\end{equation}
\end{proposition}

\begin{proof}
This is the standard Jacobi--Zariski transitivity sequence for K\"ahler
differentials~\cite{StacksDifferentials}.  The first map sends a parameter
differential from \(A\) to its image in the global algebra; the second quotient
forgets parameter motion.
\end{proof}

The first term in \cref{eq:global-transitivity} is precisely the channel through
which sample-relative pieces couple in the absolute module.  Right exactness
does not assert that the first map is injective, nor that the sequence splits.

\subsection{Jacobian presentation}

Choose polynomial equality generators
\(F=(F_1,\ldots,F_m)\) for \(I\) in variables
\(x=(\theta,S^{(1)},\ldots,S^{(N)})\).  Then
\begin{equation}
  \cR^m\xrightarrow{J_F^{\mathsf T}}\cR^{P+\sum_iK_i}
  \longrightarrow\Om_{\cR/k}\longrightarrow0.
  \label{eq:global-jacobian-presentation}
\end{equation}
The image of \(J_F^{\mathsf T}\) records differential relations
\(dF_j=0\).  Different generating families can give different presentation
matrices while presenting isomorphic modules.

\begin{remark}[Global module versus fiber]
For a real solution \(p\), tensoring
\cref{eq:global-jacobian-presentation} with the residue field at \(p\) evaluates
all coefficients.  The resulting vector-space cokernel describes a cotangent
fiber.  Equality of such fibers, or equality of their row spaces, need not
imply equality of the original ideals or equality of their real projections.
\end{remark}

\begin{remark}[Singular points]
The module \(\Om_{\cR/k}\) and the exact sequence
\cref{eq:global-transitivity} exist at singular points.  What fails there is a
naive identification of module rank with the dimension of a smooth manifold.
Our algebraic decompositions remain valid; claims based on constant-rank
linear algebra require separate hypotheses.
\end{remark}

\section{Block-angular Jacobian structure}
\label{sec:block-jacobian}

Order the variables as
\begin{equation}
  (\theta,S^{(1)},\ldots,S^{(N)})
\end{equation}
and the polynomial equations sample by sample.  Write
\begin{equation}
  C_i=\frac{\partial F^{(i)}}{\partial\theta},
  \qquad
  D_i=\frac{\partial F^{(i)}}{\partial S^{(i)}}.
\end{equation}

\begin{lemma}[Cross-sample derivatives vanish]
\label{lem:cross-derivatives}
If \(i\ne j\), then
\begin{equation}
  \frac{\partial F^{(i)}}{\partial S^{(j)}}=0.
\end{equation}
\end{lemma}

\begin{proof}
By sample locality, no indeterminate in \(S^{(j)}\) occurs in any polynomial of
\(F^{(i)}\).  A formal partial derivative with respect to an absent
indeterminate is zero.
\end{proof}

\begin{theorem}[Block-angular presentation]
\label{thm:block-jacobian}
With the order above, the global Jacobian is
\begin{equation}
J_F=
\begin{bmatrix}
C_1&D_1&0&\cdots&0\\
C_2&0&D_2&\cdots&0\\
\vdots&\vdots&\vdots&\ddots&\vdots\\
C_N&0&0&\cdots&D_N
\end{bmatrix}.
\label{eq:block-jacobian}
\end{equation}
\end{theorem}

\begin{proof}
The parameter block can be nonzero in every sample because parameters are
shared.  The local diagonal blocks contain derivatives with respect to that
sample's state.  Every off-diagonal local block is zero by
\cref{lem:cross-derivatives}.
\end{proof}

The theorem holds over the polynomial quotient ring before evaluation and over
\(\R\) after evaluation at any real point.  These are different coefficient
structures.  Gaussian elimination is immediately available for a numerical
matrix over a field.  Symbolic elimination over a multivariate quotient ring
requires additional algebra and cannot be charged as ordinary field row
reduction without justification.

\subsection{Pointwise tangent compatibility}

At a real solution \(p\), a first-order perturbation
\((\delta\theta,\delta S^{(1)},\ldots,\delta S^{(N)})\) satisfies the
linearized equations exactly when
\begin{equation}
  C_i(p)\delta\theta+D_i(p)\delta S^{(i)}=0
  \qquad (i\in[N]).
  \label{eq:linearized-local}
\end{equation}
Eliminating \(\delta S^{(i)}\) gives a constraint on
\(\delta\theta\).  This is a tangent-extension question at a known point.  It
does not ask whether a point \(p\) exists, and it does not recover the constant
terms of the original equations.

\begin{proposition}[Local linearized extension criterion]
\label{prop:local-linear-extension}
Over a field, let \(P_i^\perp\) have row space
\(\kernel D_i(p)^{\mathsf T}\).  Then
\cref{eq:linearized-local} has a solution \(\delta S^{(i)}\) if and only if
\begin{equation}
  P_i^\perp C_i(p)\delta\theta=0.
  \label{eq:local-message-criterion}
\end{equation}
\end{proposition}

\begin{proof}
The vector \(-C_i(p)\delta\theta\) belongs to the image of \(D_i(p)\) exactly
when it is annihilated by every element of the left null space of \(D_i(p)\).
Writing a basis of that null space as the rows of \(P_i^\perp\) gives
\cref{eq:local-message-criterion}.
\end{proof}

Stacking the reduced constraints yields the root matrix used by differential
CMP.  The criterion is exact for the linear system at \(p\).

\section{Dependency topology and separators}
\label{sec:topology}

Construct a primal dependency graph whose vertices are scalar variables and in
which two variables are adjacent when they occur together in an input
constraint.  A coarse tree decomposition is obtained from the bags
\begin{equation}
  B_i=\{\theta_1,\ldots,\theta_P\}\cup S^{(i)}
  \qquad(i\in[N]),
  \label{eq:sample-bags}
\end{equation}
connected through a central parameter bag
\begin{equation}
  B_0=\{\theta_1,\ldots,\theta_P\}.
\end{equation}

\begin{proposition}[Star decomposition]
\label{prop:star-decomposition}
The bags \(B_0,B_1,\ldots,B_N\), with every \(B_i\) adjacent to \(B_0\), form
a tree decomposition of the coarse sample-local dependency graph.
\end{proposition}

\begin{proof}
Every constraint belongs to one sample and therefore all of its variables lie
in the corresponding bag \(B_i\).  Each local variable occurs in only that bag.
Each parameter occurs in the center and in every sample bag, which is a
connected subtree.  Hence the edge-cover and running-intersection conditions
hold.
\end{proof}

\begin{corollary}[Width bound]
\label{cor:treewidth}
If \(K=\max_iK_i\), then
\begin{equation}
  \boxed{\tw(G)\le P+K-1.}
\end{equation}
\end{corollary}

\begin{proof}
The largest bag in \cref{eq:sample-bags} has at most \(P+K\) vertices.
\end{proof}

Increasing \(N\) adds bags but does not enlarge their shared intersection.
This fact explains why sample-wise linear algebra can scale linearly with
\(N\) once the parameter and local dimensions are fixed.

\begin{remark}[What treewidth does not prove]
A tree decomposition specifies where elimination may occur.  An exact dynamic
program also needs a closed message language whose size and update time are
controlled.  For discrete variables of bounded domain, a table can have size
exponential in the bag width and independent of the total number of bags.  For
real polynomial constraints, a projected semialgebraic set can require new
polynomials, sign conditions, and components.  The numerical dimension of a
separator does not by itself bound this descriptive growth.
\end{remark}

\begin{remark}[Uniform and parameterized readings]
The width bound supports a parameterized analysis in \(P+K\).  A claim of
uniform polynomial time additionally requires polynomial dependence on
\(P+K\), not merely tractability for each fixed width.  The cubic linear-algebra
bound later in the paper has that dependence for evaluated matrices; the exact
semialgebraic message problem does not yet have such a bound.
\end{remark}

\section{Cotangent Message Passing}
\label{sec:cmp}

This section states the exact theorem proved by the current CMP construction.
It is a separator theorem for a linearized row space.

Let \(L_A,L_B,X_A,X_S,X_B\) be vector spaces over a field \(k\), and let
\begin{equation}
  J_A:L_A\to X_A,\quad
  J_{AS}:L_A\to X_S,\quad
  J_{BS}:L_B\to X_S,\quad
  J_B:L_B\to X_B
\end{equation}
be linear maps.  A pair \((\lambda,\mu)\in L_A\oplus L_B\) generates the
block-angular row combination
\begin{equation}
  \bigl(J_A\lambda,
  J_{AS}\lambda+J_{BS}\mu,
  J_B\mu\bigr).
  \label{eq:block-row-combination}
\end{equation}

\begin{definition}[Intrinsic CMP message]
\label{def:intrinsic-message}
The message from the \(A\)-side to the separator is
\begin{equation}
  \boxed{
  \cE(A\to S):=J_{AS}(\kernel J_A)\subseteq X_S.}
  \label{eq:intrinsic-message}
\end{equation}
Equivalently,
\begin{equation}
  \cE(A\to S)=
  \{J_{AS}\lambda:\lambda\in L_A,\ J_A\lambda=0\}.
\end{equation}
\end{definition}

\begin{theorem}[Linear separator factorization]
\label{thm:linear-separator-factorization}
For \(s\in X_S\) and \(b\in X_B\), the following are equivalent:
\begin{align}
 &\exists\lambda\in L_A,\ \exists\mu\in L_B:
 J_A\lambda=0,\quad
 s=J_{AS}\lambda+J_{BS}\mu,\quad b=J_B\mu;
 \label{eq:factor-left}\\
 &\exists e\in\cE(A\to S),\ \exists\mu\in L_B:
 s=e+J_{BS}\mu,\quad b=J_B\mu.
 \label{eq:factor-right}
\end{align}
\end{theorem}

\begin{proof}
From \cref{eq:factor-left}, set \(e=J_{AS}\lambda\); the condition
\(J_A\lambda=0\) implies \(e\in\cE(A\to S)\), yielding
\cref{eq:factor-right}.  Conversely, membership
\(e\in\cE(A\to S)\) supplies \(\lambda\in\kernel J_A\) with
\(e=J_{AS}\lambda\); substitution yields \cref{eq:factor-left}.
\end{proof}

The theorem has no rank, smoothness, or invertibility hypothesis.  It is exact
for the stated linear maps.

\subsection{Matrix form and basis independence}

Suppose row vectors of \(P_A^\perp\) form a basis of the relevant left null
space.  Then
\begin{equation}
  M_{A\to S}=P_A^\perp J_{AS}
  \label{eq:matrix-message}
\end{equation}
has row space \(\cE(A\to S)\).  A different null-space basis left-multiplies
\(P_A^\perp\) by an invertible matrix and therefore gives the same message
subspace.  Redundant rows may be removed without altering that subspace.

\begin{proposition}[Message dimension]
\label{prop:message-dimension}
If \(X_S\) is finite-dimensional, then
\begin{equation}
  \dim_k\cE(A\to S)\le\dim_kX_S.
\end{equation}
\end{proposition}

\begin{proof}
The message is a linear subspace of \(X_S\).
\end{proof}

\subsection{Recursive linearized elimination}

At a fixed point, each leaf sends a basis for the differential relations that
survive elimination of its internal coordinates.  At an internal node, child
messages are stacked with the node's own rows, internal columns are eliminated,
and the result is rank-reduced on the parent separator.  Induction using
\cref{thm:linear-separator-factorization} proves that the root message is exact
for the global \emph{linearized} row space.

For dense matrices of dimension at most \(k=P+K\), conventional null-space and
rank computations use \(O(k^3)\) field operations per bag.  With \(N\) bags,
this gives
\begin{equation}
  O\bigl(N(P+K)^3\bigr)
  \label{eq:cmp-arithmetic-count}
\end{equation}
field operations for the evaluated linear problem.  The coefficient field and
the cost of representing its elements remain part of any bit-complexity claim.

\subsection{Limit of the theorem}

The message in \cref{eq:intrinsic-message} is homogeneous first-order data.  It
contains directions generated by linear combinations of derivatives.  It does
not, by definition, contain the constant terms of the input polynomials or a
description of their real zero set.  The next section tests whether it can
nevertheless serve as a complete feasibility message and gives a negative
answer.

\section{Global exactness and decision completeness}
\label{sec:global-exactness}

An exact separator algorithm for \(\ExactNNTP\) needs messages with
zeroth-order semantics.  If \(T_v\) is a rooted subtree, \(s\) its separator
assignment, and \(x\) its internal variables, the required invariant is
\begin{equation}
  s\in\mathsf{Msg}_v
  \quad\Longleftrightarrow\quad
  \exists x\;F_{T_v}(x,s).
  \label{eq:desired-semantic-invariant}
\end{equation}
Soundness is the left-to-right implication; completeness is the converse.  At
the root, nonemptiness of the message must be equivalent to global feasibility.

The semantic sets \(Q_i\) in \cref{eq:sample-projection} satisfy this invariant
for leaf subtrees.  Their conjunction satisfies it for the dataset.  What is
missing is a polynomially bounded representation and a polynomial-time
projection/merge procedure.

\subsection{A one-ReLU audit family}

Fix a positive real target \(y\).  Consider one shared scalar parameter
\(\theta\), local preactivation \(z\), and local activation \(a\), with
constraints
\begin{equation}
  z-\theta=0,\qquad a=\ReLU(z),\qquad a-y=0.
  \label{eq:toy-relu}
\end{equation}

\begin{lemma}[Exact projection of the audit family]
\label{lem:toy-projection}
For \(y>0\), the exact surviving parameter set of
\cref{eq:toy-relu} is \(Q_y=\{y\}\).
\end{lemma}

\begin{proof}
The affine equation gives \(z=\theta\) and the target equation gives \(a=y>0\).
The ReLU equation then forces \(z=a=y\), so \(\theta=y\).  Conversely,
\((\theta,z,a)=(y,y,y)\) satisfies every constraint.
\end{proof}

Use the complementarity polynomial for the ReLU equation and set
\begin{equation}
  f_1=z-\theta,\qquad
  f_2=a(a-z),\qquad
  f_3=a-y.
\end{equation}
At the satisfying point \(p_y=(\theta,z,a)=(y,y,y)\), their Jacobian rows in
the column order \((\theta,z,a)\) are
\begin{equation}
  J_y=
  \begin{bmatrix}
    -1&1&0\\
    0&-y&y\\
    0&0&1
  \end{bmatrix}.
  \label{eq:toy-jacobian}
\end{equation}
For a row combination \((\lambda_1,\lambda_2,\lambda_3)\), eliminating the
local \((z,a)\) columns imposes
\begin{equation}
  \lambda_1-y\lambda_2=0,qquad
  y\lambda_2+\lambda_3=0.
  \label{eq:toy-elimination}
\end{equation}
The surviving parameter coefficient is \(-\lambda_1\).

\begin{theorem}[Scalar cotangent message loses the target]
\label{thm:scalar-message-universal}
For every \(y\ne0\), the message
\begin{equation}
  E_y:=\{-\lambda_1:\lambda_1-y\lambda_2=0,
  \ y\lambda_2+\lambda_3=0\}
\end{equation}
equals \(\R\).
\end{theorem}

\begin{proof}
Let \(s\in\R\).  Choose
\(\lambda_1=-s\), \(\lambda_2=-s/y\), and \(\lambda_3=s\).  These choices
satisfy \cref{eq:toy-elimination} and give \(-\lambda_1=s\).  Hence
\(\R\subseteq E_y\), while the reverse inclusion is immediate.
\end{proof}

\begin{corollary}[No exact decoder from the scalar message alone]
\label{cor:no-message-decoder}
There is no function
\begin{equation}
  D:\{\text{subsets of }\R\}\times\R\to\{\mathsf{false},\mathsf{true}\}
\end{equation}
such that, for every \(y>0\) and \(\theta\in\R\),
\begin{equation}
  D(E_y,\theta)=\mathsf{true}
  \quad\Longleftrightarrow\quad \theta\in Q_y.
\end{equation}
\end{corollary}

\begin{proof}
By \cref{thm:scalar-message-universal}, \(E_1=E_2=\R\).  By
\cref{lem:toy-projection}, \(1\in Q_1\) and \(1\notin Q_2\).  Thus the same
input pair \((\R,1)\) would have to receive both Boolean values.
\end{proof}

This counterexample occurs at \(y>0\), where the ReLU is on its active linear
branch.  The failure is not caused by the switching boundary.

\subsection{The symbolic version of the obstruction}

Even before evaluation, formal differentiation erases the target constant:
\begin{equation}
  \frac{\partial}{\partial x_j}(a-y)
  =\frac{\partial}{\partial x_j}(a-y')
  \qquad\text{for every variable }x_j.
\end{equation}
Therefore the symbolic Jacobian row family of the audit architecture is
independent of \(y\) in its target row.  Yet two positive-target samples sharing
\(\theta\) are jointly feasible exactly when their targets agree.  In
particular, targets \((1,1)\) form a yes-instance and targets \((1,2)\) form a
no-instance with the same first-order target rows.

\begin{theorem}[Cotangent-only global exactness is false]
\label{thm:cotangent-only-false}
The intrinsic first-order message
\(J_{AS}(\kernel J_A)\), whether evaluated at the positive satisfying point or
formed solely from the symbolic Jacobian rows, is not a complete invariant of
exact parameter projection for all instances of \(\ExactNNTP\).
\end{theorem}

\begin{proof}
The evaluated claim follows from \cref{cor:no-message-decoder}.  The symbolic
claim follows because the family of derivatives of the target polynomial is
independent of its constant, whereas exact feasibility in
\cref{eq:toy-relu} depends on that constant.
\end{proof}

\subsection{What remains possible}

\Cref{thm:cotangent-only-false} does not rule out every algorithm that uses the
original equations together with cotangent summaries.  It proves that any
correct algorithm must preserve additional zeroth-order information.  The
tautologically complete message
\begin{equation}
  \mathsf{ZMsg}_v
  :=\{s:\exists x_{\operatorname{int}(T_v)}F_{T_v}(x,s)\}
  \label{eq:zeroth-order-message}
\end{equation}
satisfies \cref{eq:desired-semantic-invariant}.  Computing
\cref{eq:zeroth-order-message} is existential projection, however, and general
real quantifier elimination can create a much larger representation
\cite{Collins1975,Renegar1992I,BasuPollackRoy2006}.

The unresolved repair theorem is therefore precise: find a representation
class containing every message in \cref{eq:zeroth-order-message} for the stated
ReLU training systems, prove closure under leaf construction, conjunction, and
projection, and bound both representation length and exact Turing running time
by a polynomial in the original input length.  No such theorem is established
by the present first-order framework.

\section{CMP decision algorithms: proved and required forms}
\label{sec:algorithm}

Two procedures must be distinguished.  The first is an exact algorithm for an
evaluated linearized system.  The second is the desired but currently
unconstructed exact feasibility decider.

\subsection{Differential CMP}

Given numerical Jacobian blocks at a specified point \(p\), differential CMP
performs the following operations.

\begin{enumerate}
  \item For each sample, form \(C_i(p)\) and \(D_i(p)\).
  \item Compute a basis \(P_i^\perp\) for
  \(\kernel D_i(p)^{\mathsf T}\).
  \item Form the local parameter message
  \begin{equation}
    M_i=P_i^\perp C_i(p).
  \end{equation}
  \item Stack the messages and remove linearly dependent rows.
  \item Accept a proposed parameter direction \(\delta\theta\) if and only if
  the root matrix annihilates \(\delta\theta\).
\end{enumerate}

\begin{theorem}[Correctness of differential CMP]
\label{thm:differential-cmp-correct}
At a specified real point \(p\), the root test accepts
\(\delta\theta\) if and only if there exist local directions
\(\delta S^{(i)}\) satisfying all linearized equations
\begin{equation}
  C_i(p)\delta\theta+D_i(p)\delta S^{(i)}=0.
\end{equation}
\end{theorem}

\begin{proof}
Apply \cref{prop:local-linear-extension} independently to every sample and
conjoin the resulting conditions.  Equivalently, recursively apply
\cref{thm:linear-separator-factorization} along the star decomposition.
\end{proof}

This procedure assumes the point \(p\) is already supplied.  It neither tests
whether \(p\) exists nor decides the original instance.

\subsection{The required exact procedure}

An exact decision algorithm would need a finite message representation
\(R_v\) for every subtree with a denotation
\(\denote{R_v}\subseteq\R^{S_v}\).  Its operations must satisfy:

\begin{description}[leftmargin=3.4cm,style=nextline]
  \item[Leaf construction]
  \(\denote{R_v}\) is exactly the projection of the leaf's ReLU
  constraints to its separator.

  \item[Merge]
  Combining child representations denotes the intersection of their lifted
  feasible sets with the current bag constraints.

  \item[Contraction]
  Eliminating internal variables denotes existential projection, with both
  soundness and completeness.

  \item[Root test]
  The algorithm returns \textsc{yes} exactly when the root denotation is
  nonempty.

  \item[Complexity]
  Representation length and the time for every operation are polynomial in
  the original binary input length.
\end{description}

If arbitrary quantifier-free formulas over the reals are used as
representations, the first four requirements follow from standard logical
semantics.  The fifth does not follow.  If only the cotangent subspaces of
\cref{def:intrinsic-message} are used, the fifth is inexpensive but the first
and third fail by \cref{thm:cotangent-only-false}.

\begin{problem}[Polynomial semantic message]
\label{prob:polynomial-semantic-message}
Construct a representation and algorithms satisfying all five requirements
for every finite instance of \(\ExactNNTP\), or prove that no such
representation exists under an explicitly stated computational model.
\end{problem}

\subsection{No hidden root oracle}

The root test must be included in the complexity analysis.  A procedure that
propagates compact differential messages and then invokes an unrestricted
semialgebraic feasibility solver at the root has not placed the language in
\(P\); it has moved the original difficulty into the final operation.  The
same applies if a ``message entry'' is an unevaluated existential formula whose
nonemptiness remains ETR-hard.

\section{Time, space, and bit complexity}
\label{sec:complexity}

\subsection{Arithmetic complexity of differential CMP}

Let \(K=\max_iK_i\) and \(k=P+K\).  Suppose every evaluated Jacobian block is
an explicitly supplied matrix over a field and has dimensions \(O(k)\).
Classical dense elimination computes a null-space basis, multiplies the
surviving blocks, and row-reduces the result using \(O(k^3)\) field operations
per bag.  Therefore
\begin{equation}
  T_{\mathrm{arith}}(N,P,K)=O\bigl(N(P+K)^3\bigr).
  \label{eq:arithmetic-complexity}
\end{equation}
This is a valid bound for \cref{thm:differential-cmp-correct}.

If \(N,P,K\le n\), then
\begin{equation}
  N(P+K)^3
  \le n(2n)^3
  =8n^4.
  \label{eq:uniform-arithmetic-inequality}
\end{equation}
Thus the number of field operations is uniformly \(O(n^4)\) once the matrices
are present.  For fixed \(P,K\), \cref{eq:arithmetic-complexity} is linear in
\(N\).

\subsection{Space}

A streaming implementation that processes and discards leaf matrices uses
\begin{equation}
  O\bigl(V+E+(P+K)^2\bigr)
  \label{eq:streaming-space}
\end{equation}
field elements in addition to constant bookkeeping per bag, where \(V+E\)
denotes the explicitly stored constraint incidence structure.  Retaining all
messages uses
\begin{equation}
  O\bigl(V+E+N(P+K)^2\bigr)
\end{equation}
field elements.

\subsection{Why arithmetic counting is not a P theorem}

Membership in \(P\) is a bit-complexity statement.  The following data must be
specified and bounded:

\begin{enumerate}
  \item the representation of every coefficient;
  \item the bit length of intermediate matrix entries;
  \item the cost of exact zero testing and rank decisions;
  \item the representation of any algebraic real at which a Jacobian is
  evaluated; and
  \item the size and construction cost of the zeroth-order message needed for
  global feasibility.
\end{enumerate}

For an explicitly given rational matrix, fraction-free elimination controls
coefficient growth and yields polynomial bit complexity in the matrix
dimensions and input bit length~\cite{Bareiss1968}.  This fact does not close
the present proof: a feasible point can have irrational coordinates and is not
part of the input, while finding or deciding the existence of such a point is
the original nonlinear problem.

Similarly, treating symbolic entries as elements of a multivariate quotient
ring changes the operation model.  Rank, kernels, and equality testing over
that ring cannot be charged as unit-cost field operations.  Their
representations can require polynomial reduction, localization, or ideal
membership computations.

\begin{theorem}[Conditional bit-complexity implication]
\label{thm:conditional-bit}
Assume there is an exact CMP representation satisfying
\cref{prob:polynomial-semantic-message}, with total bit running time
\begin{equation}
  N\,\poly(P+K,L)
  \label{eq:desired-bit-bound}
\end{equation}
and a polynomial-time root nonemptiness test.  Then
\(\ExactNNTP\in P\).
\end{theorem}

\begin{proof}
The instance dimensions and coefficient lengths are bounded by its explicit
binary encoding length.  Under the assumed bound, the stated CMP procedure is
a deterministic polynomial-time Turing decider whose acceptance semantics is
exactly \(\Feas(I)\).
\end{proof}

\begin{remark}[Current status]
The hypothesis of \cref{thm:conditional-bit} is \openstatus.  Equation
\cref{eq:uniform-arithmetic-inequality} proves only the arithmetic inequality;
it does not construct a global decider.
\end{remark}

\section{The ETR bridge}
\label{sec:etr-reduction}

To obtain \(\ExistsR\)-hardness, a reduction must begin with an explicitly
encoded ETR sentence
\begin{equation}
  \varphi=\exists x_1,\ldots,x_m:\Psi(x_1,\ldots,x_m)
\end{equation}
and construct in polynomial time an instance \(I_\varphi\) of the exact
language in \cref{sec:problem} such that
\begin{equation}
  \varphi\text{ is true}
  \quad\Longleftrightarrow\quad
  \Feas(I_\varphi).
  \label{eq:etr-reduction-spec}
\end{equation}

Related neural-network training languages are known to be
\(\ExistsR\)-complete~\cite{AbrahamsenKleistMiltzow2021}.  A citation to that
result does not by itself prove \cref{eq:etr-reduction-spec}: the architecture,
activation function, loss/threshold convention, permitted trainable
parameters, and exact acceptance predicate must match, or a polynomial
model-conversion reduction must be given.

\subsection{Required gadget ledger}

The reduction must provide the following constructions with forward and
reverse correctness proofs.

\begin{center}
\begin{tabular}{p{0.22\linewidth}p{0.50\linewidth}p{0.14\linewidth}}
\toprule
Object & Required invariant & Status\\
\midrule
Rational constant & A designated signal/parameter equals the encoded rational. & Open\\
Signed variable & An unrestricted real variable is represented despite nonnegative ReLU outputs. & Open\\
Equality & Two represented expressions are equal iff the exact output constraints hold. & Open\\
Addition & A represented output equals the sum of represented inputs. & Open\\
Multiplication & A represented output equals the product of two represented unknowns, with legal sharing in the network model. & Open\\
Weak/strict sign & The source sign predicate is preserved using only permitted exact constraints. & Open\\
Conjunction & All source constraints hold simultaneously under one shared assignment. & Open\\
Reuse & Multiple occurrences of one source variable denote one real value without exponential copying. & Open\\
\bottomrule
\end{tabular}
\end{center}

Unknown weights multiply activations in an affine layer, so the network
equations contain bilinear terms.  Turning that observation into a multiplication
gadget requires more than noting the term's existence: both operands must be
represented as source variables, their signs must be handled, and all uses
must remain synchronized by the permitted parameter-sharing rules.

\subsection{Size and bit bounds}

For a valid many-one reduction, the construction must satisfy
\begin{equation}
  |\langle I_\varphi\rangle|
  \le \poly(|\langle\varphi\rangle|).
\end{equation}
This includes the number of nodes, samples, trainable incidences, local state
variables, and the bit lengths of rational constants.  Correctness must be
proved in both directions; a gadget that introduces unintended ReLU branches
is insufficient.

\begin{problem}[Model-matching ETR reduction]
\label{prob:etr-reduction}
Prove \cref{eq:etr-reduction-spec} for the exact target-equality language of
\cref{sec:problem}, with polynomial construction size and coefficient bit
length, or modify the language and every upstream theorem in a stated,
consistent manner.
\end{problem}

\begin{remark}[Current status]
Membership in \(\ExistsR\) is proved in \cref{prop:nntp-in-existsr}.  Hardness
for this exact formal language remains \openstatus until
\cref{prob:etr-reduction} is closed.
\end{remark}

\section{The conditional complexity collapse}
\label{sec:collapse}

The class-theoretic part of the proposed argument is short once its two
substantive certificates exist.

\begin{theorem}[Conditional \(P=\ExistsR\)]
\label{thm:conditional-p-existsr}
Assume:
\begin{enumerate}
  \item \(\ExactNNTP\in P\); and
  \item \(\mathrm{ETR}\le_p\ExactNNTP\).
\end{enumerate}
Then \(P=\ExistsR\).
\end{theorem}

\begin{proof}
The reduction and closure of \(P\) under polynomial-time many-one reductions
imply \(\mathrm{ETR}\in P\), hence \(\ExistsR\subseteq P\).  The standard
containment \(P\subseteq\ExistsR\) gives equality.
\end{proof}

\begin{corollary}[Conditional \(P=NP=\ExistsR\)]
\label{cor:conditional-collapse}
Under the hypotheses of \cref{thm:conditional-p-existsr},
\begin{equation}
  P=NP=\ExistsR.
\end{equation}
\end{corollary}

\begin{proof}
Use \(P\subseteq NP\subseteq\ExistsR=P\).
\end{proof}

\begin{remark}[No unconditional collapse claim]
The first hypothesis is not obtained from differential CMP because
cotangent-only global exactness is false by
\cref{thm:cotangent-only-false}.  The second hypothesis requires the
model-matching construction of \cref{prob:etr-reduction}.  Consequently,
\cref{thm:conditional-p-existsr,cor:conditional-collapse} are logical
implications, not proofs that their hypotheses hold.
\end{remark}

The Lean development accompanying this paper formalizes an abstract version of
these implications by packaging the two assumptions as a CMP decision
certificate and an ETR bridge certificate.  It does not instantiate either
certificate for the exact language.

\section{Globality, singularities, and ReLU switching boundaries}
\label{sec:boundary}

The exact ReLU encoding covers the switching boundary semantically.  At
\(z=a=0\), however, the complementarity polynomial has vanishing gradient.
Thus an argument that identifies cotangent rank with the codimension of a
smooth real manifold cannot be extended to this point without additional
work.

\subsection{Strata and their intersections}

For a fixed activation pattern, ReLU constraints reduce to affine equalities
and weak sign conditions.  The complete network is the union over activation
patterns, with different strata meeting where one or more preactivations
vanish.  An exact decision procedure must account for all strata and their
intersections.  Enumerating patterns directly can be exponential in the number
of gates.

Square-slack variables encode the weak inequalities without selecting a
pattern, but they do not make the resulting variety smooth.  At a zero slack,
the derivative of \(u^2\) also vanishes.  Hence polynomialization preserves the
exact feasible projection while potentially adding singular coordinates.

\subsection{Three levels of information}

The following objects must not be conflated:
\begin{enumerate}
  \item the global module \(\Om_{\cR/\R}\) over the coordinate ring;
  \item its fiber, or an evaluated Jacobian row space, at a chosen real point;
  \item the zeroth-order semialgebraic projection of the real solution set.
\end{enumerate}
The tensor and K\"ahler decompositions concern the first object.  Differential
CMP operates on the second.  Exact-NNTP feasibility is determined by the third.

\subsection{The obstruction is stronger than a boundary problem}

Although boundary completeness remains necessary for any all-ReLU theorem,
the audit family in \cref{sec:global-exactness} lies on the active branch with
\(y>0\).  Its failure shows that a smooth-locus extension cannot by itself
upgrade first-order CMP to a zeroth-order decider.  Both issues must be solved:
the message must retain exact feasible-set information, and it must handle
singular switching strata.

\section{Related work}
\label{sec:related-work}

\paragraph{The existential theory of the reals.}
ETR is the nonemptiness problem for semialgebraic sets presented by existential
formulas.  Its role as a complexity class, known complete problems, and the
containments between NP and PSPACE are surveyed by Schaefer, Cardinal, and
Miltzow~\cite{SchaeferCardinalMiltzow2024}.  General decision and quantifier
elimination algorithms include cylindrical algebraic decomposition
\cite{Collins1975}, the PSPACE method of Canny~\cite{Canny1988}, and the
complexity analyses of Renegar~\cite{Renegar1992I}; see also
\cite{BasuPollackRoy2006}.

\paragraph{Exact neural network training and verification.}
Abrahamsen, Kleist, and Miltzow proved \(\ExistsR\)-completeness for a neural
network training problem~\cite{AbrahamsenKleistMiltzow2021}.  Their result makes
precise why a claimed polynomial-time exact training decider requires an
equally precise comparison of formal languages.  ReLU network verification has
also motivated specialized satisfiability procedures such as
Reluplex~\cite{KatzEtAl2017}.  Verification of a fixed network and existential
training of its weights are related semialgebraic problems but have different
quantified variables and should not be identified without a reduction.

\paragraph{K\"ahler differentials.}
K\"ahler differentials provide a functorial algebraic representation of
derivations, conormal relations, and tangent/cotangent information.  The
Jacobian presentation, base change, tensor-product formula, and transitivity
sequence used here are standard commutative algebra
\cite{Eisenbud1995,StacksDifferentials}.  Their availability over singular
rings does not make a first-order invariant a complete description of real
points.

\paragraph{Treewidth and elimination.}
Tree decompositions expose small interfaces between otherwise separated
subproblems~\cite{Bodlaender1996}.  Exact message passing is efficient when the
message domain has a bounded representation closed under local combination and
elimination.  In this paper the star decomposition is elementary; the research
question is the representation complexity of real semialgebraic separator
messages.

\paragraph{Exact linear algebra.}
The arithmetic-operation count for CMP uses conventional dense elimination.
Fraction-free methods such as Bareiss elimination control intermediate integer
growth for exact rational matrices~\cite{Bareiss1968}.  These methods do not
replace nonlinear real projection or construct a feasible evaluation point.

\section{Discussion and repair program}
\label{sec:discussion}

The structural decomposition remains useful despite the failed completeness
claim.  It identifies exactly where samples interact and provides a compact
description of infinitesimal compatibility.  The counterexample identifies the
minimum semantic information that any exact extension must add.

\subsection{Potential message enrichments}

The following are research directions rather than established polynomial-time
repairs.

\paragraph{Semialgebraic projection messages.}
Using the exact set in \cref{eq:zeroth-order-message} gives perfect semantics.
The unresolved issue is whether this special family has polynomial descriptive
complexity and supports polynomial-time elimination.  General quantifier
elimination does not supply the desired bound.

\paragraph{Affine or homogenized first-order messages.}
Augmenting a derivative matrix with constant terms prevents the immediate loss
of \(y\) in the audit family and is exact for affine constraint systems.  For
nonlinear systems, an affine row space still does not generally determine the
zero set or its projection.

\paragraph{Higher jets.}
Higher-order derivative data can retain more of a bounded-degree polynomial.
A successful method would need a uniform order, a basis-independent merge
operation, correct treatment of inequalities and real components, and a
polynomial representation bound.  None follows from the first-order separator
theorem.

\paragraph{Ideals, real radicals, and elimination ideals.}
An algebraic message could retain an elimination ideal rather than a cotangent
subspace.  Ordinary elimination ideals describe Zariski closures of complex
projections; exact real feasibility can additionally require inequalities and
real-radical information.  Gr\"obner or quantifier-elimination representations
may grow substantially.  The special block structure would have to yield a new
bound.

\subsection{Priority order for the target theorem}

Further work should proceed in the following dependency order:
\begin{enumerate}
  \item propose a zeroth-order message representation that distinguishes the
  audit family;
  \item prove its exact semantic invariant for leaf, merge, contraction, and
  root operations;
  \item attack it with small disconnected, singular, and incompatible-target
  instances;
  \item prove a polynomial representation-size theorem;
  \item derive a Turing bit-complexity bound;
  \item formalize the model-matching ETR reduction; and
  \item only then instantiate the conditional collapse theorem.
\end{enumerate}

This ordering prevents an arithmetic speedup for an auxiliary linear problem
from being mistaken for a complexity classification of the nonlinear decision
language.

\subsection{A publishable unconditional core}

Independently of the open collapse target, the following package is
unconditional: exact ReLU polynomialization including the boundary; the
shared-base tensor decomposition; the relative K\"ahler decomposition; the
block-angular Jacobian; the intrinsic linear CMP factorization; and the
in-scope impossibility theorem for cotangent-only exact decoding.  A future
positive algorithm can build on this audited base.

\section{Conclusion}
\label{sec:conclusion}

Exact ReLU training has a strong sample-separable algebraic structure.  Its
coordinate ring is assembled from sample-local factors over one shared
parameter algebra, its relative K\"ahler differentials decompose after base
change, and its Jacobian is block-angular.  Cotangent Message Passing exploits
this structure to eliminate local coordinates exactly in the linearized row
space, using messages whose vector-space dimension is bounded by the separator.

Exact feasibility requires more.  The one-ReLU audit family proves that the
existing first-order message can be identical for problems with different
zeroth-order parameter projections.  Accordingly, the present results do not
establish a polynomial-time decider for \(\ExactNNTP\), and therefore do not
establish \(P=\ExistsR\) or \(P=NP\).  The missing theorem is now explicit: a
polynomially represented separator message whose denotation is the exact real
existential projection at every subtree, together with polynomial Turing-model
operations and a model-matching ETR reduction.  Resolving that theorem---or
proving a barrier for the proposed message families---is the next mathematical
task.


\appendix
\section{Full exact-ReLU algebra}
\label{app:relu}

This appendix expands the proof obligations behind
\cref{lem:relu-complementarity,lem:relu-polynomialization}.

\subsection{Case partition}

The equation \(a(a-z)=0\) gives the algebraic alternatives
\(a=0\) and \(a=z\).  The two weak inequalities select the valid half of each
branch:
\begin{center}
\begin{tabular}{cccc}
\toprule
Algebraic branch & Consequence of inequalities & ReLU regime & Included boundary\\
\midrule
\(a=0\) & \(z\le0\) & inactive & \(z=0\)\\
\(a=z\) & \(z\ge0\) & active & \(z=0\)\\
\bottomrule
\end{tabular}
\end{center}
At \(z=a=0\), both algebraic branches meet, but they represent the same ReLU
assignment.  The complementarity description therefore has no spurious
\((z,a)\)-point.

\subsection{Square-slack fibers}

If \(a>0\), the equation \(a=u^2\) has two real slack witnesses.  If \(a=0\),
it has the unique witness \(u=0\).  The same holds for \(a-z=v^2\).  Hence the
projection of the slack variety to \((z,a)\) is exactly the ReLU graph, although
the map is generally many-to-one.

For one gate, the equality-only encoding adds two variables and three
equalities, all of degree at most two.  For \(G\) ReLU gates and \(N\) samples,
the construction adds \(2GN\) real variables and \(3GN\) polynomial
equalities.  With an explicit architecture and dataset, this is polynomial in
the input length.

\subsection{Boundary Jacobian}

For
\begin{align}
 h_1&=a-u^2,\\
 h_2&=a-z-v^2,\\
 h_3&=a(a-z),
\end{align}
the Jacobian in \((z,a,u,v)\) at the boundary origin is
\begin{equation}
  \begin{bmatrix}
   0&1&0&0\\
   -1&1&0&0\\
   0&0&0&0
  \end{bmatrix}.
\end{equation}
The complementarity row vanishes and the slack derivatives vanish.  Thus
existential square polynomialization provides exact semantics but does not
remove singular behavior from the equality presentation.

\subsection{Formal verification correspondence}

The Lean theorem
\path{exactReLU_iff_complementarity} proves
\cref{lem:relu-complementarity}.  The theorem
\path{exactReLU_iff_polynomialized} proves
\cref{lem:relu-polynomialization}.  Separate lemmas cover active, inactive, and
boundary regimes.  These results concern the scalar gate and are composed
sample-by-sample in the mathematical encoding.

\section{Tensor products over the shared parameter algebra}
\label{app:tensor}

We give a more explicit two-factor proof of \cref{thm:quotient-tensor}.  Let
\(B_1=A[X]\), \(B_2=A[Y]\) for disjoint finite variable sets and let
\(I_1\lhd B_1\), \(I_2\lhd B_2\).  The canonical isomorphism
\begin{equation}
  \alpha:B_1\otimes_A B_2\xrightarrow{\sim}A[X,Y]
\end{equation}
sends \(f(X)\otimes g(Y)\) to \(f(X)g(Y)\).  Let \(J\) be the ideal of the
tensor product generated by the images of \(I_1\otimes B_2\) and
\(B_1\otimes I_2\).  Right exactness of tensor product gives canonical
surjections whose composite induces
\begin{equation}
 (B_1/I_1)\otimes_A(B_2/I_2)
 \xrightarrow{\sim}(B_1\otimes_AB_2)/J.
\end{equation}
Under \(\alpha\), the ideal \(J\) maps to
\(I_1A[X,Y]+I_2A[X,Y]\).  Therefore
\begin{equation}
 \frac{A[X,Y]}{I_1A[X,Y]+I_2A[X,Y]}
 \cong (B_1/I_1)\otimes_A(B_2/I_2).
\end{equation}

Associativity of tensor products supplies the finite iterated statement.  No
ordering of samples is mathematically preferred; different parenthesizations
are related by the canonical associators.

\subsection{Why base choice matters}

If one instead forms a tensor product over \(k\), then the two copies of
\(A=k[\theta]\) remain independent.  For example,
\begin{equation}
  k[\theta]\otimes_k k[\theta]
  \cong k[\theta^{(1)},\theta^{(2)}],
\end{equation}
whereas
\begin{equation}
  k[\theta]\otimes_{k[\theta]}k[\theta]
  \cong k[\theta].
\end{equation}
The latter is the correct model for a single shared trainable parameter vector.

\subsection{Real points do not commute with arbitrary projection claims}

The tensor identity is an identity of algebras.  It implies that the global
equations contain no hidden cross-sample polynomial except through the shared
base.  It does not imply
\begin{equation}
  \Proj_\theta V(I)=V(I\cap A)
\end{equation}
as equality of real sets without closure and real-algebraic qualifications.
This is why the semantic projection in \cref{prop:semantic-projection} is kept
separate from elimination ideals.

\section{K\"ahler exact sequences and base change}
\label{app:kahler}

For completeness, we record the universal-property maps used in
\cref{thm:relative-differentials}.  Given \(A\)-algebras \(B,C\) and
\(D=B\otimes_A C\), define
\begin{equation}
 \delta:D\to
 (\Om_{B/A}\otimes_BD)\oplus(\Om_{C/A}\otimes_CD)
\end{equation}
on pure tensors by
\begin{equation}
 \delta(b\otimes c)=
 (db\otimes(1\otimes c),\;dc\otimes(b\otimes1)).
\end{equation}
This is an \(A\)-derivation.  Conversely, restricting the universal derivation
of \(D\) along \(B\to D\) and \(C\to D\) gives maps from the two summands to
\(\Om_{D/A}\).  The maps are inverse because they agree with the universal
derivations on pure tensors.

Applied iteratively to \(D=\cR\), this yields
\cref{eq:relative-differential-decomposition}.  The tensor factors on the
right are essential: \(\Om_{\cR^{(i)}/A}\) is naturally a module over the local
ring, while the direct sum must be compared as a module over the global ring.

\subsection{Presentation for a quotient}

If \(P=k[x_1,\ldots,x_n]\), \(I=(f_1,\ldots,f_m)\), and \(B=P/I\), the conormal
sequence is
\begin{equation}
 I/I^2\longrightarrow
 \Om_{P/k}\otimes_PB
 \longrightarrow\Om_{B/k}\longrightarrow0.
\end{equation}
Under the basis \(dx_1,\ldots,dx_n\), the first map sends the class of \(f_i\)
to
\begin{equation}
 df_i=\sum_j\frac{\partial f_i}{\partial x_j}dx_j.
\end{equation}
Choosing the displayed generators of \(I\) therefore gives the Jacobian
presentation \cref{eq:jacobian-presentation}.  The presentation need not be
minimal, and the first arrow need not be injective.

\subsection{Evaluation}

For a real point \(p:B\to\R\), tensoring with \(\R\) through \(p\) gives
\begin{equation}
 \R^m\xrightarrow{J_F(p)^{\mathsf T}}\R^n
 \longrightarrow \Om_{B/k}\otimes_{B,p}\R\longrightarrow0.
\end{equation}
This operation retains the derivatives evaluated at \(p\) and discards
polynomial coefficients that vanish under evaluation.  It cannot serve as a
procedure for deciding whether any real point \(p\) exists.

\section{Full separator factorization proof}
\label{app:separator}

Define the global linear map
\begin{equation}
 T:L_A\oplus L_B\to X_A\oplus X_S\oplus X_B
\end{equation}
by
\begin{equation}
 T(\lambda,\mu)=
 (J_A\lambda,J_{AS}\lambda+J_{BS}\mu,J_B\mu).
\end{equation}
We wish to eliminate the \(X_A\)-coordinate from \(\im T\).  Its exact
restriction to zero in that coordinate is
\begin{align}
 \mathcal Z
 &:={\{(s,b): (0,s,b)\in\im T\}}\\
 &=\{(J_{AS}\lambda+J_{BS}\mu,J_B\mu):
        J_A\lambda=0\}.
 \label{eq:restricted-global-image}
\end{align}
Since \(J_{AS}\lambda\) ranges exactly over
\(\cE(A\to S)=J_{AS}(\kernel J_A)\), we obtain
\begin{equation}
 \mathcal Z=
 \{(e+J_{BS}\mu,J_B\mu):e\in\cE(A\to S),\ \mu\in L_B\}.
\end{equation}
This is \cref{thm:linear-separator-factorization} as an equality of sets.

\subsection{Uniqueness at the level of subspaces}

The construction \(J_{AS}(\kernel J_A)\) is canonical for the given pair of
linear maps.  A matrix \(P_A^\perp\) is only a choice of generators.  If
\(P\) and \(P'\) are full-row-rank matrices with the same row space, there is an
invertible change-of-basis matrix \(U\) with \(P'=UP\), and hence
\begin{equation}
 \rowspan(P'J_{AS})=\rowspan(UPJ_{AS})=\rowspan(PJ_{AS}).
\end{equation}
If dependent rows are allowed, equality of the image subspace still follows
from the intrinsic definition.

\subsection{Recursive invariant for linear systems}

For a rooted tree of block-linear systems, let \(L_v\) be the row-coefficient
space in the subtree and let \(X_{S_v}\) be the parent-separator coordinate
space.  The message is the image on \(X_{S_v}\) of those row combinations whose
coordinates internal to the subtree vanish.  The leaf case is
\cref{def:intrinsic-message}; the induction step is
\cref{thm:linear-separator-factorization} after stacking child messages with
the current rows.  At the root, the message therefore gives exactly the
relations in the global linear row space after all nonroot coordinates have
been eliminated.

This invariant concerns generated linear relations.  Replacing ``row
combination whose internal coordinates vanish'' by ``assignment extendable to
a nonlinear real solution'' changes the semantic object and invalidates this
proof.

\section{Global CMP exactness audit}
\label{app:cmp-exactness}

This appendix spells out the logical incompatibility behind
\cref{cor:no-message-decoder}.

For positive \(y\), let
\begin{equation}
 \mathsf{One}(y,\theta)
 :\Longleftrightarrow
 \exists z,a:\ z=\theta\land a=\ReLU(z)\land a=y.
\end{equation}
Then
\begin{equation}
 \mathsf{One}(y,\theta)\quad\Longleftrightarrow\quad\theta=y.
 \label{eq:one-relu-iff}
\end{equation}
Let \(E_y\) be the scalar message in
\cref{thm:scalar-message-universal}.  Assume for contradiction that a decoder
\(D\) satisfies
\begin{equation}
 D(E_y,\theta)\quad\Longleftrightarrow\quad\mathsf{One}(y,\theta)
 \qquad(y>0).
\end{equation}
At \((y,\theta)=(1,1)\), the right side is true, so \(D(E_1,1)\) is true.  Since
\(E_1=E_2=\R\), this implies \(D(E_2,1)\) is true.  But
\cref{eq:one-relu-iff} makes \(\mathsf{One}(2,1)\) false, a contradiction.

\subsection{Two-sample Boolean separation}

Define
\begin{equation}
 \mathsf{Two}(y_1,y_2)
 :\Longleftrightarrow
 \exists\theta:\mathsf{One}(y_1,\theta)
 \land\mathsf{One}(y_2,\theta).
\end{equation}
For positive targets,
\begin{equation}
 \mathsf{Two}(y_1,y_2)\quad\Longleftrightarrow\quad y_1=y_2.
\end{equation}
Thus \(\mathsf{Two}(1,1)\) is true and \(\mathsf{Two}(1,2)\) is false.  The
formal derivative of the target equation \(a-y\) is the same for both target
values.  A message construction that depends only on those derivatives cannot
distinguish the Boolean answers.

\subsection{Scope of the impossibility theorem}

The theorem rules out a decoder whose only semantic summary of an eliminated
subproblem is the cotangent message subspace.  It does not rule out carrying
the original target constants, nonlinear equations, sign data, elimination
ideals, or complete semialgebraic projections.  Each enrichment must be
analyzed for semantic sufficiency and polynomial representation size.

\subsection{Lean correspondence}

The formal development contains:
\begin{itemize}
  \item \texttt{scalarCMPMessage\_eq\_univ};
  \item \texttt{scalarCMPMessage\_one\_eq\_two};
  \item \texttt{exact\_projection\_one\_diff\_two};
  \item \texttt{no\_exact\_decoder\_from\_scalarCMPMessage\_alone}; and
  \item \texttt{first\_order\_data\_not\_complete\_for\_exact\_relu\_decision}.
\end{itemize}
These theorems formalize the exact in-scope obstruction rather than a generic
counterexample about unrelated varieties.

\section{Complexity ledger}
\label{app:complexity}

Let a dense local matrix have \(O(k)\) rows and columns.  Conventional Gaussian
elimination uses \(O(k^3)\) arithmetic operations and \(O(k^2)\) stored field
elements.  Matrix multiplication and rank reduction have the same classical
cubic upper bound.  Over \(N\) sample bags this yields
\(O(Nk^3)\) operations and, when streamed, \(O(k^2)\) working field elements.

For \(k=P+K\) and \(N,P,K\le n\), monotonicity gives
\begin{align}
 N(P+K)^3
 &\le n(n+n)^3\\
 &=n(2n)^3\\
 &=8n^4.
\end{align}
This inequality has been machine-checked in the Lean development.

\subsection{Conditional rational-matrix bit bound}

Suppose an \(r\times c\) matrix with rational entries of bit length at most
\(L\) is explicitly given.  Clearing denominators produces an integer matrix
with polynomially related bit length.  Fraction-free elimination represents
intermediate entries through minors; determinant bounds then give polynomial
bit lengths in \(r,c,L\).  This supports a polynomial Turing bound for the
linear-algebra phase~\cite{Bareiss1968}.

The premises are not automatically satisfied by symbolic CMP:
\begin{itemize}
  \item a satisfying point is not given;
  \item its coordinates need not be rational;
  \item symbolic entries lie in a polynomial or quotient ring; and
  \item the zeroth-order message may contain several semialgebraic components.
\end{itemize}

\subsection{Checklist for a future bit theorem}

A complete proof of \cref{eq:desired-bit-bound} must state:
\begin{enumerate}
  \item a canonical binary encoding for every message;
  \item polynomial bounds on message length after every merge and projection;
  \item exact algorithms for equality, sign, and emptiness in that encoding;
  \item coefficient-height bounds under every operation;
  \item the number of operations along the tree; and
  \item the cost of the root acceptance test.
\end{enumerate}
Without these items, an \(O(n^4)\) count of abstract ring or field operations is
not a proof of membership in \(P\).

\section{ETR reduction obligations}
\label{app:etr}

This appendix records a proof template for
\cref{prob:etr-reduction}.  Normalize the source formula to a conjunction of
primitive constraints only after proving that the normalization has polynomial
size and preserves satisfiability.  A suitable primitive source language might
contain
\begin{equation}
 x_i=c,\qquad x_i+x_j=x_k,\qquad x_ix_j=x_k,
 \qquad x_i\ge0,
\end{equation}
but the completeness and normalization theorem for the chosen language must be
cited or proved.

For each primitive constraint, a network gadget must come with a relation
\(R_G\) between its interface values and satisfy
\begin{equation}
 R_G(\text{interface})
 \quad\Longleftrightarrow\quad
 \text{source primitive holds}.
\end{equation}
The reverse implication is essential: additional trainable weights or
activation branches must not create spurious satisfying assignments.

\subsection{Composition requirements}

Gadgets must compose while preserving one value for every repeated source
variable.  If source variables are represented by shared trainable parameters,
the reduction must show that all required occurrences of a parameter are legal
under the architecture encoding.  If they are represented by activations, the
reduction must supply equality/copy gadgets and handle negative values.

The final construction must use exact target equalities as specified in
\cref{sec:problem}; importing a theorem for total squared error below a
threshold requires an explicit conversion.  Likewise, a theorem for a
different activation or for trainable architecture cannot be cited as hardness
for this language without a reduction.

\subsection{Audit table for every gadget}

\begin{center}
\begin{tabular}{>{\raggedright\arraybackslash}p{0.28\linewidth}>{\raggedright\arraybackslash}p{0.60\linewidth}}
\toprule
Audit item & Required proof\\
\midrule
Forward correctness & Every source assignment extends to a network witness.\\
Reverse correctness & Every network witness restricts to a valid source assignment.\\
Boundary behavior & Zero preactivations introduce no extra source solutions.\\
Size & Nodes, edges, samples, and states grow polynomially.\\
Height & All rational coefficients have polynomial bit length.\\
Uniformity & A Turing machine constructs the encoded instance in polynomial time.\\
Scope & The output is exactly an instance of \(\ExactNNTP\).\\
\bottomrule
\end{tabular}
\end{center}

Only after every row is discharged may the manuscript state
\begin{equation}
  \mathrm{ETR}\le_p\ExactNNTP
\end{equation}
as an unconditional theorem.

\section{Lean 4 verification notes}
\label{app:lean}

The accompanying Lean 4 development is on the isolated repository branch
\texttt{cmp-lean-verification}.  This manuscript was initialized from commit
\texttt{fcca060331defca7226db3e0aecc8bc3c93fcbdf}.  GitHub Actions run 74 for
that commit completed successfully.  The workflow runs \texttt{lake build} and
then rejects occurrences of \texttt{sorry}, \texttt{admit}, or a line beginning
with \texttt{axiom} in the \texttt{CMPLean} directory.

\begin{center}
\small
\begin{tabular}{>{\raggedright\arraybackslash}p{0.31\linewidth}p{0.58\linewidth}}
\toprule
Lean file & Verified content used by this manuscript\\
\midrule
\texttt{ReLU.lean} & Complementarity, square-slack polynomialization, and active/inactive/boundary cases.\\
\texttt{DecisionProblem.lean} & Abstract Boolean existential feasibility semantics with shared parameters.\\
\texttt{ExactProjection.lean} & Feasibility iff existence in the shared-parameter projection.\\
\texttt{KahlerCore.lean} & Transitivity exactness and relative-map surjectivity from Mathlib.\\
\texttt{BlockPolynomial.lean} & Cross-sample formal partial derivatives vanish.\\
\texttt{AmbientTensor.lean} & Ambient shared-base polynomial tensor equivalence.\\
\texttt{Separator.lean} & Basis-free linear separator factorization.\\
\texttt{MessageSubmodule.lean} & Message as image of a kernel and separator-dimension bound.\\
\texttt{ComplexityArithmetic.lean} & The inequality \(N(P+K)^3\le8n^4\) under dimension bounds.\\
\texttt{ReLUDecisionTests.lean} & One- and two-sample positive-target feasibility.\\
\texttt{MessageCompleteness}\newline\texttt{Audit.lean} & Identical scalar CMP messages with different exact projections; no message-only decoder.\\
\texttt{ZeroOrderAdversarial.lean} & Target constants vanish from first derivatives while Boolean feasibility changes.\\
\texttt{CollapseLogic.lean} & Conditional class-collapse logic from abstract decision and reduction certificates.\\
\bottomrule
\end{tabular}
\end{center}

The formalization does not construct the CMP decision certificate or the ETR
bridge certificate.  In particular, the theorem named
\texttt{P\_eq\_NP\_eq\_ExistsR} accepts those certificates as hypotheses.  Its
successful compilation verifies the final logical implication, not the missing
global exactness or reduction premises.

The current formal audit is unusually informative because it includes a
negative theorem internal to the exact ReLU scope.  Any repaired CMP definition
should first be instantiated on the one-ReLU and incompatible-target tests
before a new global certificate is attempted.


\bibliographystyle{alpha}
\bibliography{references}

\end{document}
